\documentclass[11pt,letterpaper]{article}
\usepackage[margin=0.87in,headheight=15pt]{geometry}
\usepackage{fontspec}
\setmainfont{Latin Modern Roman}
\setsansfont{Latin Modern Sans}
\setmonofont{DejaVu Sans Mono}[Scale=0.79]
\usepackage{unicode-math}
\setmathfont{Latin Modern Math}
\usepackage{microtype,booktabs,longtable,array,enumitem,fancyhdr,fvextra,xcolor}
\usepackage{xurl}
\usepackage[colorlinks=true,linkcolor=teal!65!black,urlcolor=teal!65!black]{hyperref}
\definecolor{ink}{HTML}{183449}
\definecolor{pale}{HTML}{EEF3F6}
\setlength{\parskip}{5pt}
\setlength{\parindent}{0pt}
\setlength{\emergencystretch}{1em}
\setcounter{tocdepth}{1}
\newcolumntype{P}[1]{>{\raggedright\arraybackslash}p{#1}}
\setlist{nosep,leftmargin=1.5em,topsep=5pt}
\pagestyle{fancy}
\fancyhf{}
\fancyhead[L]{\small\sffamily ProveIt / Lean proof simplification}
\fancyhead[R]{\small\sffamily 26 October 2026}
\fancyfoot[C]{\small\thepage}
\renewcommand{\headrulewidth}{0.3pt}
\newcommand{\code}[1]{\nolinkurl{#1}}
\newcommand{\src}[2]{\href{https://github.com/VladimirReshetnikov/ProveIt/blob/60e9c5d620c6c387d932876cfa2d9376f64ea56c/Computability/DiophantineConsistency/Lean/DiophantineConsistency/#1.lean}{\nolinkurl{#1.lean}} (lines #2)}
\newcommand{\snapshot}{\code{60e9c5d620c6c387d932876cfa2d9376f64ea56c}}
\DefineVerbatimEnvironment{Lean}{Verbatim}{fontsize=\small,breaklines=true,breakanywhere=true,samepage=true,frame=single,rulecolor=\color{black!15},framesep=5pt}
\hypersetup{pdftitle={Shortening and simplifying the DiophantineConsistency Lean proofs},pdfauthor={Prepared for Vladimir Reshetnikov}}
\begin{document}
\begin{titlepage}
\vspace*{0.5in}
{\sffamily\color{ink}\Large ProveIt research engineering review\par}
\vspace{0.28in}
{\sffamily\bfseries\color{ink}\Huge Shorter, clearer Lean proofs\par}
\vspace{0.12in}
{\sffamily\bfseries\color{ink}\LARGE DiophantineConsistency\par}
\vspace{0.28in}
{\Large Concrete refactorings, checked prototypes, and an adoption plan\par}
\vspace{0.35in}
Prepared for Vladimir Reshetnikov\\
26 October 2026
\vspace{0.35in}

\textbf{Source snapshot.} \snapshot\par
\textbf{Pinned environment.} Lean 4.32.0; Mathlib revision
\code{81a5d257c8e410db227a6665ed08f64fea08e997}.

\vspace{0.25in}
\colorbox{pale}{\parbox{0.94\textwidth}{\vspace{7pt}
\textbf{Main recommendation.} Concentrate on repeated semantic interfaces:
variable-layout transport, checked list folds, residual aggregation,
finite certificate assembly, and source-register decoding. The repository
already contains many of the generic mathematical theorems needed for
this work. Making their clients use them consistently is the largest
opportunity.\par
\vspace{5pt}
The accompanying archive contains a precise 40-line proof-only patch,
standalone Lean prototypes, source measurements, and reproducible checks.
Component reductions are distinguished from unmeasured whole-project
savings.\vspace{7pt}}}
\vfill
\textbf{Scope.} This is a proposal and prototype review of existing proofs.
It does not claim a new consistency theorem, a completed integration of the
current polynomial with its machine interpretation, or a measured global
compilation speedup. Existing repository sources were kept unchanged.
\end{titlepage}

\tableofcontents
\newpage
\section{Findings and priorities}

The best route to shorter proofs is to make the repeated mathematical
interfaces explicit, then retain the original executable objects and prove
their correctness through those interfaces. In this project, proof length,
the size of a represented program, the size of a polynomial, and the cost of
kernel checking are separate quantities. A successful proof refactor can
leave all program and polynomial data definitionally unchanged.

There are three especially promising levels of work.
First, small proof-only changes can immediately remove local induction and
transport boilerplate. Second, a few reusable semantic APIs can replace
the same list, witness, and sum-of-squares arguments in several modules.
Third, the generated circuit proof files can use the certificate machinery
already proved in the repository, substantially reducing their materialized
proof source while retaining named semantic intermediate results.

\begin{longtable}{@{}P{0.06\textwidth}P{0.40\textwidth}P{0.30\textwidth}P{0.16\textwidth}@{}}
\toprule
Order & Proposal & Evidence or deliverable & Effort\\\midrule
\endfirsthead
\toprule Order & Proposal & Evidence or deliverable & Effort\\\midrule
\endhead
1 & Replace seven list-lookup inductions & Exact patch; 40 physical lines saved & Small\\
2 & Share witness-layout semantics & Compact constructors with expression-equality checks & Small\\
3 & Share checked-history folds and scan recurrences & Six generic laws, checked in Lean & Small--medium\\
4 & Centralize executable sum-of-squares semantics & Three repeated evaluation proofs; shared degree/cost API & Small--medium\\
5 & Use parent power certificates in generated circuits & 193 power declarations; 1,538-line old component & Medium\\
6 & Specialize local DAG equations on demand & 999 existing per-index wrappers & Medium\\
7 & Replace half-radix case splits by an exact identity & Stronger generic arithmetic lemmas, checked & Small\\
8 & Abstract finite chunk assembly & Repeated decoder and rail dispatch; retain checked leaves & Medium\\
9 & Share physical-field residues and source facts & Literal duplicate residues and repeated Pell row decoding & Medium\\
10 & Derive carry recurrence from a smaller certificate & Existing digit bound plus exact decomposition suffice & Medium\\
11 & Factor object-language proof transport & 27 completeness calls; reusable sealed-validity bridge & Medium\\
12 & Move projection laws and generic gate proofs earlier & Existing APIs blocked by dependency direction & Medium\\
\bottomrule
\end{longtable}

The order reflects risk and ease of review, rather than a claim that the
earliest item yields the largest line reduction. The power certificate work
offers the largest concrete component compression; the witness and scanner
interfaces offer the strongest improvement to day-to-day proof maintenance.
Several further ideas are developed below: a dynamic-seed primitive-recursive
iterator, Boolean reflection, a small copied-field interface, and disciplined
proof profiling.

\section{Reviewed snapshot and measurements}

The review uses the repository's \code{main} commit \snapshot. The checkout
was pinned before source inspection. The active library rooted at
\code{Computability/DiophantineConsistency/Lean} contains 546 Lean files
and 73,971 physical lines. Of these, a filename/path classification identifies
220 check or audit files and 326 other files, including the aggregate import
module. The latter group contains 66,857 physical lines. The whole requested
directory contains 711 tracked Lean files; 165 lie in historical candidate
or provenance material outside the active library. They are not included
in the active-library line totals or treated as independent refactoring
targets.

\begin{center}
\begin{tabular}{@{}lr@{}}
\toprule
Active-library measurement & Value\\\midrule
Lean files & 546\\
Physical source lines & 73,971\\
Nonempty code lines after masking comments/strings & 59,832\\
Occurrences of \code{theorem} in code & 4,970\\
Occurrences of \code{omega} / \code{nlinarith} & 1,681 / 563\\
Occurrences of \code{ring} / \code{decide} & 686 / 4,080\\
\code{maxHeartbeats} / \code{maxRecDepth} settings & 134 / 177\\
\bottomrule
\end{tabular}
\end{center}

These are lexical measurements, not kernel proof-term sizes or timings.
The comment-aware scan found no \code{sorry}, \code{admit},
\code{native_decide}, or explicit axiom declarations in the active source.
That observation does not replace transitive axiom audits or compilation.
The supplied \code{scripts/inventory.py}, CSV files, and JSON inventory
make the counting method and file hashes inspectable.

The largest proof-source targets are:

\begin{center}
\begin{tabular}{@{}p{0.58\textwidth}rr@{}}
\toprule
Module & Physical lines & Code lines\\\midrule
\code{SplitCircuitSteps59.lean} & 3,445 & 2,604\\
\code{GapCircuitSteps66.lean} & 3,122 & 2,375\\
\code{SplitCircuit59.lean} & 1,286 & 1,223\\
\code{SplitJointCircuit59.lean} & 1,252 & 1,189\\
\code{GapCircuit66.lean} & 1,226 & 1,163\\
\code{ShortSystem.lean} & 1,117 & 1,036\\
\bottomrule
\end{tabular}
\end{center}

The raw circuit tables and the correctness proofs serving them have
different roles. Literal data should remain reviewable. The primary target
is the repetitive proof layer around that data.

\section{Immediate proof-only improvement: use Lean's list API}

Seven local proofs repeat induction on a list to connect optional indexing
with membership. Lean 4.32.0 already supplies both directions, with the
required types:

\begin{Lean}
theorem lookup_mem {α : Type} {xs : List α} {i : Nat} {a : α}
    (h : xs[i]? = some a) : a ∈ xs :=
  List.mem_of_getElem? h

theorem exists_lookup {α : Type} (xs : List α) (a : α)
    (h : a ∈ xs) : ∃ i : Nat, xs[i]? = some a :=
  List.getElem?_of_mem h
\end{Lean}

\begin{center}
\begin{tabular}{@{}p{0.51\textwidth}rr@{}}
\toprule
Module & Proofs & Lines removed\\\midrule
\code{AnnotatedRename} & 1 & 5\\
\code{AnnotatedBound} & 1 & 5\\
\code{AnnotatedProofChecker} & 1 & 8\\
\code{AnnotatedAssumption} & 2 & 11\\
\code{AnnotatedRuleSide} & 2 & 11\\\midrule
Total & 7 & 40\\
\bottomrule
\end{tabular}
\end{center}

The exact sites are \src{AnnotatedRename}{110--117},
\src{AnnotatedBound}{50--57}, \src{AnnotatedProofChecker}{177--187},
\src{AnnotatedAssumption}{40--57}, and \src{AnnotatedRuleSide}{18--35}.
The supplied \code{patches/lookup-api.patch} replaces proof bodies only.
All declaration names, statements, and numeric definitions remain identical.
Its two polymorphic replacement proofs compile under the pinned compiler,
and \code{git apply --check} succeeds against the snapshot. The complete
patched project modules were not replayed in this review; that distinction
is recorded in the validation evidence.

This is an excellent first change because its scope is small and its
benefit is exactly measurable. It also illustrates a broader search
method: inspect locally reproved structural lemmas before designing new
automation. Existing library facts often remove an entire proof instead
of merely shortening its tactic script.

\section{Witness layouts: prove the transport once}

\subsection{The repeated argument}

\src{Representation}{155--256} repeatedly proves that renaming input and
witness coordinates commutes with evaluation. In particular,
\code{Rep.conjunction} and \code{Rep.disjunction} each introduce the same
left and right evaluation equalities, using function extensionality and
case analysis. \code{reindex} and \code{existsBlock} contain closely
related arguments. A private one-sided helper also exists in
\src{QuarticCompiler}{95--98}; the proposal is to make the general
semantic law available at the representation layer.

The central law is:
\begin{Lean}
natAssignment v w ∘ Sum.map f g =
  natAssignment (v ∘ f) (w ∘ g)
\end{Lean}
Its proof is just \code{funext i; cases i <;> rfl}. Together with the
existing \code{Expr.eval_rename}, it eliminates the repeated left/right
transport proofs.

The second law is the existential equivalence between a single concatenated
witness block and its two components:
\begin{Lean}
(∃ w : Fin (m + n) → Nat,
  P (w ∘ Fin.castAdd n) (w ∘ Fin.natAdd m)) ↔
    ∃ a b, P a b
\end{Lean}
The forward map projects the two blocks. The reverse map uses
\code{Fin.addCases}. This works when either block has length zero, so
special empty-witness cases are unnecessary.

\subsection{Preserve the executable data}

The supplied \code{WitnessLayout.lean} defines compact conjunction and
disjunction constructors. Their witness counts are unchanged and their
expression fields are definitionally equal to the originals; the file
contains explicit \code{rfl} checks of expression equality. Thus the
proposal simplifies correctness proofs without changing the represented
polynomials, variable order, or number of witnesses.

The original two constructor declarations occupy 59 physical lines;
the compact versions occupy 32. This comparison excludes the shared
helper API. Its one-time cost must be included in a net repository
comparison, and it should be amortized across \code{reindex},
\code{existsBlock}, and further layout clients. The immediate gain is
conceptual: arithmetic correctness is no longer interrupted by repeated
function-coordinate bookkeeping.

\subsection{Move existing projection lemmas to their natural home}

The desired witness-count simplification lemmas already exist in
\src{RepresentationSize}{136--147}, with graph analogues at lines
193--209. However, \src{DirectComposition}{59--68} still unfolds whole
constructors to compute counts. It cannot import the size module without
reversing the present dependency direction.

Move the small \code{rfl} projection lemmas next to the constructors or
into a lightweight module that imports only the constructors. The size
module can then concentrate on global bounds. This changes which facts
are available early, rather than adding duplicate facts. It also prevents
the simplifier from examining large expression and proof fields merely
to calculate a natural-number projection.

\section{Checked histories and numeric scanners}

\subsection{One invariant-preserving history fold}

The list processors in \code{AnnotatedRename}, \code{AnnotatedBound},
and \code{AnnotatedProofChecker} have the same recursive shape: validate
the next line against a history; prepend the emitted result; continue;
or reject. Rename and bound rows emit themselves; proof lines emit a
judgment. The appropriate abstraction therefore has two types and two
functions:

\begin{Lean}
def run (valid : Line → List Judgment → Bool)
    (emit : Line → Judgment) :
    List Line → List Judgment → Option (List Judgment)
  | [], seen => some seen
  | line :: lines, seen =>
      if valid line seen then
        run valid emit lines (emit line :: seen)
      else none
\end{Lean}

The generic soundness proof assumes exactly the local invariant step:
if all existing judgments satisfy \(P\), and a line passes validation,
then its emitted judgment satisfies \(P\). Induction on the remaining
lines proves the invariant for the final history.

The duplicated clusters are \src{AnnotatedRename}{196--235},
\src{AnnotatedBound}{136--175}, and
\src{AnnotatedProofChecker}{206--250}. They contain \code{run_append},
\code{run_sound}, \code{Produces.append}, and \code{Produces.step}.
The supplied \code{ValidatedFold.lean} proves the shared versions.
The existing \code{Produces} interface is useful; its implementation and
laws should become shared while public module-local names remain thin
specializations. Formula-specific soundness remains in the current row
or line soundness theorem.

A conservative migration keeps the old executable \code{run} definitions
and proves their equality with the common fold. A later migration may
replace those definitions by transparent specializations, but only after
checking any syntax-sensitive size claims. Extensional equality of
programs does not imply equality of their encoded abstract syntax trees.

\subsection{Separate recurrence from encoding}

Five scanner modules independently establish the same iterator equations:
\code{AnnotatedRenameScan}, \code{AnnotatedBoundScan},
\code{AnnotatedLineScan}, \code{AnnotatedAxiomContext}, and
\code{AnnotatedContext}. Their \code{scan_succ} and fixed-terminal-state
proofs depend on iteration, not on pairing codes or formulas.

For an arbitrary state type, define
\begin{Lean}
def scan (step : State → State) (fuel : Nat)
    (initial : State) : State :=
  fuel.rec initial (fun _ acc => step acc)
\end{Lean}
Then prove once:
\begin{Lean}
scan step (fuel + 1) initial = scan step fuel (step initial)
step fixed = fixed → scan step fuel fixed = fixed
\end{Lean}
The two supplied recurrence proofs compile without any axioms.
They can replace the repeated induction proofs while preserving the old
numeric definitions through transparent unfolding.

The next stage is to share the simulation between encoded lists and typed
folds. Concrete targets are \src{AnnotatedRenameScan}{71--87},
\src{AnnotatedBoundScan}{67--83}, \src{AnnotatedLineScan}{132--152},
and \src{AnnotatedAxiomContext}{130--151}. The common simulation should
take a \code{step_cons} equation, an acceptance test, a history update,
and terminal/rejection encodings.

Two details are essential. Ordinary histories use zero as the empty code,
whereas an empty quoted context uses one. Also, correctness for raw codes
must continue to cover malformed inputs and arbitrary supplied fuel.
Proving only the behavior of well-typed encodings would weaken that contract.

\subsection{A better primitive-recursive iterator interface}

\src{DirectScalarArithmetic}{103--107} already provides
\code{iterate}, with a fixed numeric seed. The useful extension is
\code{iterateFrom}, accepting realized seed and fuel functions, plus a
parameterized step that can still access the original input. The repeated
constructions occur in \src{AnnotatedContextPR}{47--54},
\src{AnnotatedBoundScan}{144--158}, \src{AnnotatedLineScan}{231--245},
and \src{AnnotatedRenameScanPR}{48--75}.

The result should remain an explicit \code{R f} object with supplied
program data. A bare existence theorem \code{Nat.Primrec f} would lose
the data needed by the compiler and its size accounting. The expensive
validator is deliberately used once in
\src{AnnotatedLineScan}{201--226}; preserve that sharing. Require
equality of the constructed \code{.val} program, or recheck the exact
node-count theorems, before claiming unchanged encoded size.

\subsection{Share tag dispatch and Boolean reflection}

\src{AnnotatedRenameLocal}{63--87} and
\src{AnnotatedBoundLocal}{55--79} repeat the same tag partition:
\(0\)--\(1\), \(2\), \(3\)--\(5\), \(6\)--\(7\), and rejection for
tags at least \(8\). A single arithmetic dispatch lemma can certify that
the threshold implementation matches this partition, parameterized by
the four branch results.

The scanner interface can also be strengthened from an acceptance
equivalence to an exact value equation. Combine
\code{check = 1 ↔ valid = true} and
\code{check = 0 ∨ check = 1} into \code{check = truth valid}.
This removes repeated true/false reasoning in the two \code{step_cons}
proofs. The range premise is essential: the acceptance equivalence alone
allows a rejecting numeric result of two.

\section{Residual aggregation and generic gate semantics}

\subsection{One executable sum-of-squares API}

Three modules contain equivalent private
\code{eval_sumExpressions} proofs:
\src{DAGQuartic}{83--86}, \src{NaturalDAGQuartic}{110--113}, and
\src{QuadraticRegisterSystem}{70--73}. Each then proves that a finite
sum of squared residuals is zero exactly when each residual is zero,
using the same \code{List.ofFn}, finite-sum conversion, and ordered-ring
lemma. Their degree proofs also repeat the same sum and product bounds.

Introduce a public expression constructor and three reusable theorems:
\begin{Lean}
def sumSquares (rows : Fin n → Expr σ) : Expr σ :=
  QuarticSyntax.sumExpressions
    (List.ofFn fun i => rows i * rows i)

(sumSquares rows).eval v = 0 ↔
  ∀ i, (rows i).eval v = 0

(∀ i, (rows i).toPolynomial.totalDegree ≤ d) →
  (sumSquares rows).toPolynomial.totalDegree ≤ 2 * d

(∀ e ∈ es, e.binaryOps ≤ b) →
  (QuarticSyntax.sumExpressions es).binaryOps ≤
    (b + 1) * es.length
\end{Lean}

The first equivalence is over integers, or another suitable ordered
structure where a sum of squares vanishes termwise. It is not valid over
an arbitrary commutative ring. The supplied prototype deliberately uses
integer evaluation. The degree statement is an upper bound, so it does
not claim an exact degree in the presence of cancellation.

The cost generalization also absorbs the existing special case with
row cost five in \src{QuarticCost}{60--65}, and the general argument
repeated in \code{DAGQuartic}. Because the expression uses the same
ordered list of squared residuals as before, clients can retain the exact
syntax rather than replacing it by an algebraically equivalent polynomial.

Signed witness reconstruction and natural witness reconstruction need
different hypotheses. Share residual aggregation first; a fully generic
signed/natural encoding framework should be attempted only if its
instantiation overhead is smaller than the duplication it replaces.

\subsection{The semiring theorem already subsumes the natural induction}

\src{QuarticCompiler}{101--167} proves natural gate completeness and
soundness by syntax induction. \src{GatesSemiring}{72--145} proves the
corresponding results over an arbitrary commutative semiring, covering
the natural case. The two expression interpretations are already
connected by \code{Expr.eval_eq_zero_iff_splitNatural} and
\code{NatExpr.evalSemiring_nat}.

The useful refactor is architectural. Extract gate data, wire indexing,
roots, residual syntax, and value definitions into a lightweight module.
Prove generic gate semantics above it. Then derive the natural-number
results through the interpretation bridge, retaining their public names.
Currently \code{GatesSemiring} imports \code{QuarticCompiler}; simply
calling the generic theorem from the compiler would introduce a cycle.
Moving the data is the prerequisite.

Likewise, \src{QuadraticRegisterSystem}{18--33} implements an
unweighted degree and its polynomial bound, while
\src{WeightedExpressionDegree}{37--62} already proves the generic
weighted result. Constant weight one supplies the old bound. Keep an
abbreviation or compatibility equality if downstream proofs depend on
the old recursive equations. The polynomial-renaming theorem in
\code{QuarticCompiler.Gates} belongs in the expression API for the same
reason: basic syntax laws should be available without importing a compiler.

\section{Generated circuit proofs: use the existing certificates}

\subsection{Where the source volume comes from}

\code{GapCircuitSteps66.lean} contains 484 local-step theorems and 86
power theorems. \code{SplitCircuitSteps59.lean} contains 515 step theorems
and 107 power theorems. Together they occupy 6,567 physical lines.
Every local-step theorem already calls the generic
\code{ArithmeticDAG.run_node_equation}; that theorem is not missing.
The repository also already has \code{run_eq_of_node_equations},
\code{value_eq_of_node_equations}, \code{run_power}, and a certified
power-sharing rewrite. The opportunity is to adopt these existing
interfaces consistently at the generated clients.

\subsection{Replace 193 power chains by parent exponent certificates}

The old power sections occupy 685 and 853 lines respectively, for a
combined 1,538 lines. Each chain follows multiplication nodes and proves
an exponent identity. The existing \code{run_power} theorem in
\src{ArithmeticDAGPowers}{26--46} abstracts exactly this argument.

Assign a natural exponent label \(\ell(i)\) to each node. Label zero
means that the node is outside the certified pure-power fragment.
The distinguished base has label one. Each labeled non-base node must
be a multiplication of earlier labeled nodes, with
\[
\ell(i)=\ell(a)+\ell(b).
\]
The literal node table is finite, so these local conditions are checked
by \code{decide +kernel}. Topology and the existing theorem then yield
\[
\operatorname{value}(i)=\operatorname{value}(\mathrm{base})^{\ell(i)}.
\]

The supplied \code{GapPowerCertificate.lean} and
\code{SplitPowerCertificate.lean} keep the original
\code{naturalNodes} exactly. They are proof certificates for those
circuits, not optimized replacement circuits. Their initial source
representation occupies 117 lines including label tables and examples,
compared with the 1,538-line old power component. This is a component
comparison; it is not a measured integrated project reduction. Existing
named callers must be migrated or retained as thin compatibility wrappers.
The supplied metrics record the concrete versions included in the archive.

For Split, all 515 entries of the proposed label table agree with the
old-label table already present in
\src{SplitJointPower59}{64--78}. Move that parent certificate to a small
shared module rather than importing an unrelated child circuit solely to
reuse its labels. The Python generator checks labels against the actual
multiplication nodes as a diagnostic. The Lean certificate is the proof.

Do not infer a compilation speedup from the shorter source. One finite
certificate can still be expensive to reduce. Compare actual replay time,
proof object size, and peak memory before choosing its final chunking.

\subsection{Specialize low-level node equations at use sites}

The supplied \code{CircuitDispatch.lean} wraps the generic node equation
in one \code{step p x i hi} theorem. At a concrete index, its type reduces
to the same scalar equation as \code{natural_step_i}. A short example
uses six selected node equations to prove the base identity
\(p_0(x_0+1)^2\).

This creates an alternative to eagerly declaring all 999 low-level
wrappers. However, replacing every short theorem name by a verbose
specialization can make clients harder to read. The practical compromise
is a small local elaboration helper or generator that emits applications
of the existing theorem, together with named results at meaningful
register boundaries.

Do not make unrestricted node expansion a global simplification rule.
Unfolding all shared registers recursively can turn a DAG into a large
expression tree. The named base, kernel, mask, rank, and source bindings
in \code{GapCircuitRegisters66} and \code{GapCircuitRecipe66} are useful
proof boundaries. Keep those boundaries, and reduce the repetitive
low-level transport beneath them.

\section{Native arithmetic: exact identities and shared geometry}

\subsection{A stronger half-radix lemma}

For all natural \(a,b\) and positive \(s\),
\[
\boxed{\quad
(sa)\bmod(2s)+\left\lfloor\frac{sb}{2s}\right\rfloor
=s(a\bmod2)+\left\lfloor\frac b2\right\rfloor.
\quad}
\]
The proof cancels a common positive factor in the modulus and quotient.
The Lean proof uses \code{Nat.mul_mod_mul_left} and
\code{Nat.mul_div_mul_left}, after commuting \(2s\).

Since \(a\bmod2\leq1\), the expression is at most \(s+1\) whenever
\(b\leq3\), with no bound on \(a\). If \(s\) is even, the expression's
parity equals \(\lfloor b/2\rfloor\) under the same bound on \(b\).
Without that bound, the parity is instead
\(\lfloor b/2\rfloor\bmod2\).

This replaces the four-case bound proof in
\src{NativeCyclicRotation}{100--116} and the sixteen combinations used
for parity at lines 160--165. The original proof bodies occupy 15 and
four nonblank lines; proposed wrappers are one-line applications of the
generic lemmas. Their helper cost must still be counted, so this is a
local simplification and strengthening, not automatically a net line
saving across the whole project.

\code{NativeRefactorCore.lean} checks the identity, bound, and parity
lemmas using only Lean's core arithmetic infrastructure. The supplied
\code{NativeRotationWrappers.lean} shows their intended applications to
the actual \code{shiftDigits}; its integration status is stated in the
validation ledger. The already generic Boolean rotation API in
\code{NativeBooleanRotation} remains useful and does not need replacement.

\subsection{Share physical-field residue calculations}

The addition/subtraction residue proofs in
\src{CurrentClause2011}{27--39} are repeated privately in
\src{CurrentReflectedGeometry2011}{74--86}. A small
\code{CyclicFieldArithmetic} module should own these facts for total
length \(NL\) and field width \(L\). For an aligned base \(b\), field
offset \(e<L\), and arbitrary natural shift \(t\), a useful normal form is
\[
(b+e-t).\operatorname{val}\bmod L
=\bigl(L+e-(t\bmod L)\bigr)\bmod L.
\]
The positional subtraction on the left is in \(\mathbb Z/(NL)\mathbb Z\).
The bounded natural subtraction on the right is safe because
\(t\bmod L<L\). This is already the substance of
\src{CurrentCellTranslation2011}{84--94}; extract it downward instead
of importing that later module backward and creating a cycle.

The clause and marker convolution proofs both calculate the residues
of shifts 18 and 2020, then use support exclusion. A shared residue
normal form and support-zero helper would replace repeated arithmetic
with one instantiation and a small numerical check. The origin-exclusion
step in the corresponding mask lemmas can also be shared.

A particularly direct reuse is available for reflected coefficients.
\src{CurrentClause2011}{122--130} proves the residue modulo 2048.
\src{CurrentEarlyPadding2011}{95--107} reproves the corresponding
modulo-16 result by cases. After sharing the identical selector-unit
definition, project the stronger congruence using
\code{.of_dvd (by decide : 16 ∣ 2048)}.

Preserve arbitrary shifts, wraparound, and \(N=1\). A later equivalence
between aligned cells and \code{ZMod N} could simplify reindexing, but
the immediate extraction does not require a new coordinate representation.

\subsection{Decode source-register equations once}

Several Pell arguments reconstruct the same raw arithmetic equations:
\[
q=1+J\beta,\qquad L=q^2-1,\qquad
N=q(1+L)=q^3,\qquad X=qw.
\]
Concrete repetitions occur in \code{source_Y_ge_two},
\code{source_rank_eq_target}, and \code{source_scale_eq} in
\src{NativePellBootstrap}{95--115, 415--436, 498--518}, and in
\src{NativePellPower}{127--143}. Main and auxiliary classifications
also decode related norm rows separately.

Extend the existing source-register helper approach to a small set of
named arithmetic facts proved once from the literal rows. State the
primitive identities over integers and derive natural-number corollaries
with explicit nonnegativity premises where subtraction is involved.
Pell arguments then consume semantic equations, with fewer row unfoldings
and casts.

Keep raw row equations, positivity, port bounds, Pell classification, and
rank recovery in separate interfaces. Bundling all of them together can
silently add assumptions to clients. The existing \code{PortBox} already
packages port inequalities; a differently named replacement adds no value.

\section{Carry certificates and finite data assembly}

\subsection{Derive recurrence instead of storing it twice}

The existing \code{CurrentPolynomialWindow2011.Facts} record stores a
digit bound, a quotient recurrence, and an exact local decomposition.
The quotient recurrence follows from the other two. In natural arithmetic,
\[
u=d+rk,\quad d<r
\quad\Longrightarrow\quad
k=\lfloor u/r\rfloor,\quad d=u\bmod r.
\]
The strict digit bound already implies \(r>0\). No extra positivity
assumption is needed in the generic implication.

The checked \code{CarryCertificate} prototype stores carry bounds,
digit bounds, and decomposition, indexed by the existing arrays and
predecessor function. Its \code{recurrence} theorem is a derived method.
In the current window, take
\(u_i=\operatorname{convolution}_i+
\operatorname{rotation}_i+k_{i-1}\).

This reduces repeated argument lists in
\src{CurrentPhaseParity2011}{92--163} and the parallel fields in
\src{CurrentPolynomialWindow2011}{27--65}. It also makes explicit which
facts are independent. Construct the record using the existing
normalization theorem obtained from the actual polynomial. A convenient
record must not become a new unproved assumption in place of that
construction.

Preserve minimal low-level statements. Some carry-reset lemmas need only
decomposition and carry bounds, without a digit bound or recurrence.
Keep those lemmas general and offer record-based wrappers to larger
clients. Changing every theorem to demand a large certificate would
shorten notation while strengthening hypotheses unnecessarily.

\subsection{Keep checked finite chunks; share their assembly}

\src{WangTiles942Data}{210--321} proves
\code{action_eq_decoder} by repeatedly splitting an index into successive
16-entry ranges. Its rail-cache counterpart repeats the pattern at
lines 439--512. Each branch reconstructs a local index, proves bounds,
and transports an already checked chunk equality.

The generic dispatcher in \code{NativeRefactorCore.lean} uses Euclidean
division below \(Bq\), and subtraction in the final tail:
\[
i=B\lfloor i/B\rfloor+(i\bmod B),\qquad
i<Bq+r.
\]
It combines \(q\) complete blocks of size \(B\) and a final block of
size \(r\). The action table uses \((B,q,r)=(16,26,14)\); the rail
table uses \((16,13,8)\). A small finite case selection can then choose
the already authenticated chunk theorem. An estimated 90--120 lines
could be removed across these two dispatchers, after helper overhead;
this is an estimate awaiting an integrated patch.

The same principle applies to the hand-built sum assembly in
\src{CurrentKernel2011}{156--419}. Keep the manageable independently
checked chunks, but describe their combination as data and apply one
proved aggregation theorem. A list of partial totals is enough to
generate a balanced assembly tree if reduction cost warrants it.

Do not replace kernel-checked leaves with \code{native_decide}. Native
evaluation adds a compiler-dependent axiom to the proof's trust boundary,
as documented by the official Lean reference. Nor should one assume that
a single enormous \code{decide +kernel} call is cheaper than the present
chunks. Compression should remove repeated assembly logic while retaining
the ability to control kernel reduction size.

\subsection{One copied-field bound interface}

\src{CurrentCellLabels2011}{66--77, 283--299} reconstructs the same
four copied-field bounds, using offsets and widths
\((978,8),(960,9),(969,8),(942,9)\). The convolution-word proof repeats
the corresponding four block instantiations.

A small first step is \code{copiedFields_bounded}, returning the four
facts once. If more clients emerge, use a finite field descriptor with
offset, width, and packed shift, plus a generic field evaluation/bound
theorem. Preserve the actual order: packed fields are west, north, east,
south, whereas the \code{ColorsBounded} conclusion orders west, east,
north, south. A broad tuple refactor must handle that difference explicitly.

\section{Object-language proofs: centralize the semantic boundary}

\subsection{A shared bridge from arbitrary models to sealed derivations}

There are 27 occurrences of \code{completeness_inf_context} in 18
\code{Object*.lean} files. Twenty-five follow the direct pattern of
extracting a model's ZF axioms, proving semantic validity, and applying
universal closure. Examples include
\src{ObjectNaturalExpression}{181--202},
\src{ObjectMultiplication}{198--215}, and
\src{ObjectAdditionLaws}{183--206}.

The following is the intended interface, not an assertion that this
particular integration theorem has been replayed:
\begin{Lean}
theorem zfc_proves_seal_of_zf_valid (φ : Form)
    (hφ : ∀ (V : Type) (mem : V → V → Prop),
      ZFAxioms mem → ∀ ee : Nat → V, Sat mem ee φ) :
    BProv ZFCax_s [] (sealF φ)
\end{Lean}

Its proof owns completeness, the extraction of axioms from a ZFC model,
and \code{seal_valid}. Each arithmetic client supplies only its
semantic validity theorem. The private model-extraction proof in
\src{ObjectArithmetic}{143--164} and public copy in
\src{ObjectMultiplication}{169--183} can be consolidated at a low
dependency layer. A later upstream helper already exists, but importing
the much heavier bounded-consistency layer into early arithmetic merely
to save a few lines would be a poor dependency tradeoff.

Keep the quantification over arbitrary models visible. Their internal
natural numbers may be nonstandard. Replacing their arithmetic by Lean's
ordinary \code{Nat} would change the theorem. The existing
\code{ObjectOmegaSemiring}, \code{ObjectNaturalOrderedSemiring}, and
\code{ObjectNaturalSemiring.evalRel_iff} already provide the relevant
internal algebra; adding semiring instances is not a new recommendation.

Likewise, \code{ObjectNaturalExpression.evalRel_unique} currently has
no natural-input guard. A typed semiring specialization with stronger
input hypotheses does not subsume that theorem. Preserve its current
generality separately.

\subsection{Generalize closed-sentence transport across arity}

\src{ObjectQuarticCompiler}{100--116} defines a local \code{closed}
helper for the source arity and then repeats the proof for the target
arity. Quantify over an arbitrary arity \(k\), with argument
\code{(f : Expr (Fin k))}.
The target becomes another application of the same theorem.

Preferably expose the result as
\code{ObjectPolynomial.sat_noZeroForm_at_iff}, parallel to the existing
environment-specific consistency lemmas in \code{ObjectConsistency}
and \code{ObjectCodeSeal}. This is a narrow reusable interface for
closed formulas; it does not require redesigning the object language.

\section{Adoption plan and criteria}

\subsection{Stage 1: proof-only edits and shared laws}

Start with the lookup patch. Next introduce witness-layout transport,
sum-of-squares aggregation, and generic scan recurrences. Preserve old
declaration names and executable definitions, adding compatibility lemmas
where needed. Check the expression and program data before spending time
on shorter tactic scripts: \code{rfl} equality is the strongest convenient
result when it is available.

Each review should report old and new source lines including the shared
helper, the affected statements, and the axiom set of the public endpoints.
Do not count a common helper as free. Some APIs pay back only after their
second or third client migrates.

\subsection{Stage 2: reuse certificates and reduce generated wrappers}

Move parent exponent certificates into reusable modules and replace power
chains in the old generated proofs. Retain thin aliases during migration
if that keeps downstream changes small. Then introduce scoped node
specialization beneath the named semantic register lemmas. Measure source
bytes as well as lines, because dense certificate tables make physical
line counts sensitive to formatting.

For finite tables, keep the checked leaf size initially and change only
the dispatcher or sum assembly. Rechunking and changing the mathematical
representation should be separate changes, so a performance regression
has an identifiable cause.

\subsection{Stage 3: reorganize dependencies and object transport}

Move basic constructor projection laws down the import graph. Split gate
data from generic semantics before replacing natural syntax inductions.
Extract shared cyclic-field arithmetic, source-register equations, and
the arbitrary-model bridge into appropriately early modules. Then migrate
clients incrementally, preserving minimal assumptions at low levels.

Avoid importing the aggregate \code{DiophantineConsistency} module into
a small proof that needs only one or two leaf APIs. This is an engineering
recommendation based on the aggregate's broad import list, not a measured
claim about any particular build's bottleneck.

\subsection{Profile a concrete bottleneck rather than increasing budgets}

The source contains many heartbeat and recursion-depth settings. Those
are candidate locations to inspect, not proof that those files are slow.
For a representative theorem, record elaboration time, kernel-checking
cost where available, peak memory, and generated object size under the
same compiler, cache state, and concurrency. Use the pinned compiler's
profiling tools and narrow the measurement to a declaration or module.

Prefer stable semantic rewrites and a small \code{simp only} vocabulary
at a large interface. During development, use search or suggested
simplifications to discover the right facts, then keep the resulting
proof explicit. Global simplification rules that unfold full circuits,
entire source-row systems, or encoded syntax can erase the very sharing
that makes the proof tractable.

\section{Validation, limitations, and reproducibility}

Six supplied Lean files passed with Lean 4.32.0. The successful audited
declarations use only \code{propext}, \code{Classical.choice}, and
\code{Quot.sound}; several use fewer, and the generic scanner recurrence
lemmas use none. No successful audit contains \code{sorryAx} or a
native-evaluation axiom.

\begin{longtable}{@{}P{0.31\textwidth}P{0.59\textwidth}@{}}
\toprule
File & What was checked\\\midrule
\endfirsthead
\toprule File & What was checked\\\midrule
\endhead
\code{LookupAPI.lean} & Two generic replacements for list-lookup induction,
using the pinned core library.\\
\code{ValidatedFold.lean} & Checked-history append and soundness, fragment
composition, and two generic iterator laws.\\
\code{NativeRefactorCore.lean} & Exact half-radix arithmetic, bounds/parity,
finite chunk dispatch, quotient/remainder consequences, and carry recurrence.\\
\code{WitnessLayout.lean} & Compact constructors importing the original
\code{Representation}; expression fields agree with the originals by
\code{rfl}.\\
\code{SplitPowerCertificate.lean} & Original compiled \code{SplitCircuit59};
107 local power labels, the uniform power theorem, and exponent-1201 example.\\
\code{GapExactDataFixture.lean} & Exact original 484-node data, topology,
86 power labels, uniform power theorem, exponent-2200 example, and scoped
node dispatch/base identity.\\
\bottomrule
\end{longtable}

The original \code{GapCircuit66} module was terminated with exit code 137
during the focused dependency build. The original \code{SplitCircuit59}
module and the generic DAG semantic/power modules completed successfully.
For Gap, the fallback fixture copies the complete \code{naturalNodes}
definition byte-for-byte: 8,078 definition bytes, with extraction offsets
and SHA-256 hashes recorded in \code{evidence/gap_fixture_source.json}.
Only the length, topology, power and dispatch obligations needed by the
prototype are replayed there. This is a narrower validation than building
the full original module.

The production-import Gap power and dispatch candidates therefore have
exact-data fixture validation, not a successful replay through the full
original Gap import. \code{CompilerCommon.lean},
\code{NativeRotationWrappers.lean}, and
\code{SemiringSpecialization.lean} remain integration proposals without
full module replay. They are kept in a separate \code{proposals/} folder.
The exact lookup patch passes applicability checking; the five entire
patched project modules have not been compiled.

The final checked witness code supplies the existential-layout predicate
explicitly to rewriting. An initially implicit higher-order rewrite did
not match. Similarly, the dispatch prototype needs scoped node-table,
\code{Node.evalAt}, list-index and \code{Fin.mk_zero} reductions, while
leaving the evaluator itself opaque. These are concrete elaboration
details resolved in the successful prototypes, not reasons to introduce
unrestricted global simplification.

All per-file logs and the machine-readable validation manifest are included.
Reported wall times are environment-specific observations; they do not
establish a performance improvement over the old proofs.


\subsection{What the checks establish}

An exit-zero Lean compilation checks the supplied declaration's proof
in its stated environment. The expression-equality examples in the
witness prototype additionally check the proposed constructors' literal
expression data. The power prototypes, when marked passed in the ledger,
use the original circuit definitions and validate their exponent-label
conditions in Lean. Python's parallel label checks are diagnostics and
do not enlarge the trusted proof base.

The lookup patch is distinct: the two generic proof terms compile and
the patch applies cleanly, but the five entire modified modules have
not been replayed. The other large refactorings are source-grounded
designs or standalone theorem prototypes. No whole-project shortened
variant has been built, and no whole-project speedup or reduction
percentage is claimed.

\subsection{Logical boundaries to retain}

Several source facts constrain any simplification:
\begin{itemize}
\item \code{NativeFieldCycles2011.offset19_excludes_510} excludes value
510. The genuine 511 self-loop satisfies its assumptions and is tested
in the existing checks. A replacement claiming that every nonzero cycle
is impossible would be false.
\item The offset-20 exclusion applies to an arbitrary function; it must
not gain a finiteness or permutation hypothesis merely to share a wrapper.
\item Cell-translation facts include wrapping, arbitrary large shifts,
and the one-cell case. Keep these concrete regression checks.
\item The maintained current polynomial's arithmetic window and the
pending copy/label/machine composition are different proof obligations.
Refactoring arithmetic interfaces does not complete the latter.
\item Ordinary host-level arithmetic, arbitrary-model internal naturals,
and syntactic ZFC derivability remain separate semantic layers.
\end{itemize}

These are specific properties of the reviewed source, not generic
restrictions on proof automation. They identify where a superficially
shorter proof could inadvertently weaken the intended result.

\subsection{Archive contents and replay}

The archive includes this \code{report.tex}, its compiled PDF, the
Lean prototypes, the proof-only patch, source inventories and dependency
edges, label-certificate generation and checking scripts, and validation
logs. \code{README.txt} explains the exact commands and distinguishes
verified modules from integration sketches.

To reproduce the source inventory, run the supplied Python script against
a checkout at the pinned commit. To reproduce Lean checks, use the
repository's \code{lean-toolchain} and \code{lake-manifest.json}; the
replay script builds only the imports needed by the selected prototypes.
The report can be rebuilt with XeLaTeX twice. All source links in the
PDF are pinned to the reviewed commit, so later repository changes do
not silently change the evidence.

\section*{Primary sources}
\addcontentsline{toc}{section}{Primary sources}

The principal evidence is the source linked throughout this report at
\snapshot. The repository's \href{https://github.com/VladimirReshetnikov/ProveIt/blob/60e9c5d620c6c387d932876cfa2d9376f64ea56c/lean-toolchain}{\code{lean-toolchain}}
and \href{https://github.com/VladimirReshetnikov/ProveIt/blob/60e9c5d620c6c387d932876cfa2d9376f64ea56c/lake-manifest.json}{\code{lake-manifest.json}}
identify the replay environment. The
\href{https://github.com/VladimirReshetnikov/ProveIt/blob/60e9c5d620c6c387d932876cfa2d9376f64ea56c/Computability/DiophantineConsistency/STATUS.md}{project status record}
and the comments of the current endpoint modules delimit the still-open
composition work.

For the trust-boundary distinction between kernel reduction and native
evaluation, see the official
\href{https://lean-lang.org/doc/reference/latest/Axioms/}{Lean reference, Axioms},
and \href{https://lean-lang.org/doc/reference/latest/Tactic-Proofs/Tactic-Reference/}{Tactic Reference}.
Those online pages track the current manual; all executable evidence
recorded as successful here was checked against Lean 4.32.0 as recorded
in the ledger.

\end{document}
